\documentclass[12pt]{article}
\usepackage[margin=1.1in]{geometry}
\usepackage{times}
\usepackage{hyperref}
\usepackage{xcolor}
\usepackage{listings}
\lstset{basicstyle=\ttfamily\small,breaklines=true,columns=fullflexible,frame=single,
  rulecolor=\color{gray!40},backgroundcolor=\color{gray!5}}
\hypersetup{colorlinks=true,linkcolor=teal,urlcolor=teal}
\setlength{\parskip}{7pt}\setlength{\parindent}{0pt}
\linespread{1.12}

\title{\vspace{-2em}\textbf{PatentPincer}\\[2pt]
\large An autonomous patentability analyst over SerpApi prior art}
\author{Dipankar Sarkar\\ \small \texttt{github.com/doom2quake}}
\date{}

\begin{document}
\maketitle

\begin{abstract}
A prior-art search is the first thing a founder needs and the last thing they
can afford: a patent attorney charges four figures and takes a week to say
whether an idea is already patented. PatentPincer collapses that first search
into a minute. Describe an invention in one sentence and it searches the prior
art, fetches the numbered claims of the closest references, compares the
invention element by element, and returns a decision-ready verdict (CLEAR,
NEEDS\_REVIEW, HIGH\_OVERLAP, or SEARCH\_UNAVAILABLE) with the closest reference
and a next step. SerpApi's Google Patents, Patents Details and Scholar engines
are the load-bearing data backbone; the analyst is a skill graph on the
agent-core, with API spend held behind a guardrail and the verdict routed to the
right destination. The design commitment is that the system refuses rather than
guesses: a search that did not succeed, a reference set that came back empty, or
claim text that was never retrieved all fail closed, and never to a clearance.
\end{abstract}

\section{Problem}
Most founders skip the prior-art search and find out the expensive way, after
they have built. The information exists (patents are public, papers are indexed)
but assembling it, reading it against a specific claim, and turning it into a
file-or-narrow-or-walk decision is exactly the slow, costly step. The evidence is
also split: a claim can be anticipated by a granted patent or by non-patent
literature, and a search that reads only one of the two misses half the art. We
want that first pass to be fast, cited, and honest about how close the nearest
reference sits.

\section{Design}
PatentPincer runs the workflow the way an attorney actually works, as named
skills assembled into an ADK supervisor that delegates in order:
\begin{lstlisting}
1 search-prior-art  (Patents + Scholar)
2 fetch claim text  (Patents Details)
3 assess-novelty    (element-by-element)
4 render-verdict    (CLEAR|NEEDS_REVIEW|
                     HIGH_OVERLAP|
                     SEARCH_UNAVAILABLE)
\end{lstlisting}
The searcher varies its phrasing across a few Google Patents queries to widen
recall and adds one Scholar query for non-patent literature. It then calls the
Google Patents Details engine for the shortlist, because a search snippet is not
a claim and an element-by-element comparison against snippets is guessing. The
analyst decomposes the invention into claim elements and scores each element
against every numbered claim of every reference; a cell records the reference,
the claim number, the terms that matched and the excerpt they matched in, so
every judgement is checkable and overrulable. A single reference disclosing at
least 80\% of the elements anticipates (HIGH\_OVERLAP); at least 40\% disclosed
jointly is a combination risk (NEEDS\_REVIEW). The brief begins with a
machine-readable \texttt{VERDICT:} line so the outcome can be routed.

\section{Refusing rather than guessing}
The failure that matters in this product is a confident clearance built on
nothing, so four rules are enforced in the pipeline and pinned by tests. If no
Google Patents search succeeded (guardrail suppression, timeout, bad key, quota)
the verdict is SEARCH\_UNAVAILABLE and an alert is raised. If searches succeeded
but returned zero references, the verdict is at best NEEDS\_REVIEW: an empty
reference set scores zero of $N$ elements disclosed, which would otherwise read
as the strongest possible CLEAR off the weakest possible evidence. If references
were found but no claim text came back, there is nothing to compare and the run
refuses. And if the references came from the offline demo corpus rather than a
live SerpApi response, a computed CLEAR is downgraded to NEEDS\_REVIEW, because
fixture data can warn you off an idea but cannot clear one. Provenance is carried
on every record and every call (\texttt{ok}, \texttt{fixture},
\texttt{fixture\_unavailable}, \texttt{suppressed}, \texttt{error}); only
\texttt{ok} is evidence.

\section{The hero: SerpApi as the load-bearing backbone}
The product is only as good as its prior-art coverage, and that coverage is
SerpApi. Google Patents returns structured, current patent records (assignee,
priority date, links) that are impractical to scrape reliably; the Patents
Details engine returns the numbered claims the comparison actually reads; and
Scholar catches the literature that anticipates just as many claims. Remove
SerpApi and there is no evidence left to reason over, which is precisely why the
pipeline refuses to emit a verdict when a SerpApi call did not succeed. Spend is
bounded rather than trusted: every call passes an \texttt{ActionLimiter} check
before it leaves the process, and the per-cycle cap is sized to exactly one
assessment (query variants, one Scholar query, and the claim fetches), so a
runaway loop cannot burn credits. A budget that starves the claim fetch would
silently degrade the analysis, so the optional LLM narrative is given its own
spend cycle and cannot take the analyst's. The three tools are also exposed over
MCP with an explicit \texttt{run\_id} argument, since a separate process cannot
see the caller's run context. The verdict is then routed on a strict parse of the
first line: HIGH\_OVERLAP and SEARCH\_UNAVAILABLE raise an alert and open a
ticket; NEEDS\_REVIEW and CLEAR open a ticket; an unparseable brief goes to a
human rather than being guessed at.

\section{Results}
The test suite (\texttt{pytest}, 38 tests) runs with no network and no API key
and covers each rule the product depends on. Provenance: offline records are
labelled \texttt{fixture} and never \texttt{ok}, queries outside the demo corpus
are refused instead of answered with someone else's patents, the \texttt{before}
and \texttt{assignee} filters really apply, an unparseable cutoff is an error,
and a live response carrying its own \texttt{status} field cannot relabel a call.
Claim text: the Details engine yields numbered claims and matrix cells quote them
verbatim. Fail-closed: dry-run suppression, timeout, an API-level error body,
zero results, and references with no claim text each refuse rather than clear,
and a fixture CLEAR is downgraded while a live CLEAR is not. Router: a CLEAR
brief that contains the phrase "high overlap" still routes as a clearance, and a
malformed verdict line routes to a human. Spend and context: concurrent
assessments keep separate ledgers and budgets, and an explicit \texttt{run\_id}
beats the ambient context. The live SerpApi branch is exercised through a faked
\texttt{urlopen} covering success, error bodies, timeouts, empty results and
missing claims.

\section{Related work and limitations}
Automated prior-art retrieval and patentability scoring are established; the
contribution here is the end-to-end analyst path (structured SerpApi retrieval,
claim-text retrieval, element-by-element comparison, a routed verdict) packaged
so it runs in a minute, with explicit provenance and refusal semantics. The
comparison itself is deliberately plain: light suffix stemming, stopword removal
and term containment of an element in a claim, a first-pass screen rather than a
trained model or a legal opinion. Its value is that it shows its work, so a
reviewer can see which words carried a match and overrule the cell. We report no
live-key measurements: this build had no funded SerpApi key and no GCP project,
so no recorded run carries LIVE provenance and no LLM narrative was generated;
the live path is covered by fakes, which validates request construction, parsing
and labelling but not real retrieval quality. Coverage is in any case bounded by
what Patents and Scholar index, and PatentPincer informs a filing decision rather
than replacing counsel.

\section{Conclusion}
The first prior-art pass does not have to cost a week and four figures.
PatentPincer turns a one-sentence invention into a cited, decision-ready verdict
by treating SerpApi as the evidence and the agent as the analyst that reasons
over it.

\end{document}
